\scanpage{41}
\mspage{22}
de $\xi',\xi_1$ ou de l'une de ces quantités seulement
\[
(75)\qquad\begin{aligned}
\beta'&=\xi_1\sigma_\infty(\xi')^{-1}=\deleted{\varepsilon_0(\lambda\xi_1)}\xi_1(\sigma_\infty\varepsilon_0)(\lambda\xi_1)^{-1}\\
&=\lambda^{-1}\cdot\varepsilon_0^{-1}(\xi')\sigma_\infty(\xi')^{-1}.
\end{aligned}
\]
de \deleted{la \unsure{dernière}} \inserted{ces} expressions de $\beta'$ on tire
\[
(75')\qquad\begin{aligned}
\sigma_\infty(\beta')&=\sigma_\infty(\xi_1)\xi'^{-1}=\sigma_\infty(\xi_1)\varepsilon_0(\lambda\xi_1)^{-1}\\
&=\sigma_\infty(\lambda)^{-1}\underbrace{\overbrace{\rho(\xi')\xi'^{-1}}^{\alpha'^{-1}}}_{\unsure{\nearrow}\quad\deleted{\unsure{\alpha'}}}
\end{aligned}
\]
et la troisième expression de $\beta'$ \deleted{donne}

\deleted{$(\unsure{77})$ Ceci dit, l'équation \unsure{que} $(72\,b)$ \unsure{équivalent}:}
\[
\deleted{(77\qquad\beta'\sigma_\infty(\beta')=1}
\]
\deleted{soit, $\sigma_\infty(\beta')$ en posant $\beta_1=\sigma_\infty(\beta')$ \unsure{pour}}
s'écrit aussi
\[
(77)\qquad\sigma_\infty(\beta')^{-1}=\alpha'\sigma_\infty(\lambda).
\]
Ceci posé, l'équation $(72\,b)$ équivaut à
\[
(\unsure{78})\qquad\beta'\sigma_\infty(\beta')=1\qquad\text{ou encore}\quad\unsure{\sigma}_\infty(\beta')^{-1}=\beta'
\]
ce qui est invariant par la substitution $\beta'\mapsto\sigma_\infty(\beta')$ \unsure{ou|\deleted{dans}} $\beta'\mapsto\beta'^{-1}$, donc posant
\inserted{\unsure{sans|(sans} supposer \unsure{vérifiée} $(*)$)}
\[
\beta_{\unsure{1}}=\sigma_\infty(\beta')^{-1}=\alpha'\sigma_\infty(\lambda)
\]
l'équation (78) équivaut à l'équation
\[
(**)\qquad\beta_2\sigma_\infty(\beta_1)=1
\]
et implique que $\beta_1=\sigma_\infty(\beta_1)^{-1}$ donc $\beta_1=\beta'$.

En résumé, on peut tout de suite poser
\[
(\unsure{79})\qquad\beta'=\alpha'\sigma_\infty(\lambda)\qquad\text{où}\quad\alpha'=\xi'\rho(\xi')^{-1}
\]
et la relation à imposer : $(\lambda,\xi')$ est, au choix,
\margin{40}
